\bookchapter{第四十章：需求之工——化念為 AI 可行之令}

\chapterintro{若上下文之工所解者「使 AI 知其背景」，則需求之工所解者「使 AI 知此次所作何事」。此 AI 編程至顯人我之差之節也——同器同模，有者能作令人驚豔之項目，有者徒得平庸之玩。其異往往不於技之能，而於述需之能。}

\booksection{一、自晦之念至結構化之需求}

每項目皆始於一晦之念——「吾欲作一記帳之 App」「吾欲作一個人網站」「吾欲作一待辦清單之器」。此念自無過，然爾徑擲與 AI，則產之質，全繫乎運。需求之工之之，即化此晦之念為 AI 能精行之結構化之需求也。

可喜者：結構化之需求，毋須爾作數十頁之 PRD。惟須爾答五問，吾儕謂之「需求五要」：

\begin{center}
\begin{tabularx}{\linewidth}{|X|X|X|}
\hline
\textbf{其要} & \textbf{其問} & \textbf{其例：記帳 App} \tabularnewline
\hline
Who & 誰用之？ & 自用之，記日用之收支 \tabularnewline
What & 作何事？ & 一能記每筆之收支、按月而總、可視之應用 \tabularnewline
Why & 何為而作？ & 欲知己消費之構，得其不必之支出 \tabularnewline
How & 大致何以為？ & 前端用 React，數據存瀏覽器之本地，姑不作後端與登入 \tabularnewline
Constraint & 有何限制？ & 一人為之，2 週而畢，毋須上線部署 \tabularnewline
\hline
\end{tabularx}
\end{center}

\begin{center}\small\color{gray}需求五要與其例\end{center}

此五要毋須作正式之文，即於腦中疾過一遍亦足矣。要於養「先思審而後言」之習。爾將見，能明答此五問，則爾述需自結構化矣。

\booksection{二、需求拆解分層之藝}

項目之大者，一述需不足也。須拆大需為小而可獨驗之子需。此程謂之「需求之拆解」，其核之則：每子需求皆當能獨成、獨驗、獨付。

以「作一個人記帳 App」為例，合理之拆解如次：

\begin{enumerate}
  \item 成基礎收支記錄之功（能增、觀、刪其錄）；
  \item 增分類之功（予每筆之錄以飲食、交通等籤）；
  \item 增月度彙總之功能（按月統計收支之況）；
  \item 增資料可視化（以圖表展示消費結構）。
\end{enumerate}

凡每一步，皆獨立可運行可驗證之功能增量也。

需求拆解之金律：每步既畢，項目皆當為"可用之物"——縱功能至陋。慎勿使項目滯於"事事半而無一可用"之境。

此拆解之法有二利。其一，使君恆能見進展——每成一步，君之 App 即多一功能，此正回饋於長期專案至要（人皆需成就感也）。其二，使 AI 每次任務皆足夠小，以降「AI 理解偏差」之風險——任務愈小，AI 理解愈準確，產出之質愈高。

\booksection{三、PRD 撰寫模板與實踐}

較大之專案（如一有前後端、多功能模組之全棧應用），僅憑需求五要素或不足。此時君須撰一份更正式之專案需求文檔（PRD）。一易誤解之點：PRD 非寫予 AI 觀——乃寫予君自己，助君理清思路，而後君方能予 AI 更精確之指令。

一份善 PRD 不必甚長，然當含六部：專案概述（一段話說清所為）、功能列表（每功能點為何）、技術方案（用何技術棧、大致架構）、資料模型（核心資料結構）、界面描述（主要頁面之貌）、開發計畫（分幾步而成）。

\begin{lstlisting}
# 個人記賬應用 PRD
## 項目概述
一輕量個人記賬之應用，備收支記錄、分類、月度匯總與數據可視化。
## 功能列表
1. 記錄收支：金額、類型(收入/支出)、分類、備註、日期
2. 分類管理：預設餐飲、交通、購物、娛樂、其他
3. 月度彙總：按月計總收入、總支出、結餘
4. 消費可視化：餅圖展示各分類佔比
5. 數據導出：導出爲CSV
## 技術方案
- 前端：React + TypeScript + Tailwind CSS
- 數據存儲：localStorage（無需後端）
- 圖表：使用Chart.js
- 構建：Vite
## 數據模型
Record {
  id: string
  amount: number
  type: 'income' | 'expense'
  category: string
  note: string
  date: string  // YYYY-MM-DD
}
## 開發計劃
Step 1: 基礎記錄功能 (CRUD)
Step 2: 分類管理
Step 3: 月度彙總
Step 4: 數據可視化
Step 5: 數據導出
\end{lstlisting}

注意此 PRD 之撰寫風格：用簡潔之列表與結構化格式，而非長篇大論。每功能點皆有明確定義，資料模型以簡單型別描述。此風格之 PRD 不惟君自覽清楚，君予之於 AI 時，AI 亦能精確理解每需求之點。

\booksection{四、實戰：從想法到方案之完整推演}

徒紙上談兵無味，吾儕以一完整之例遍走需求工程之全程。設君之念曰：「吾欲作一學習追蹤器，能追吾每日所學者何」。

第一步，遍核需求五要素。

\begin{itemize}
  \item Who：個人使用。
  \item What：一器也，能記每日學習之內容、追學習之時長、顧學習之進度。
  \item Why：欲養持續學習之習，見己之所積。
  \item How：先作網頁版，資料存於瀏覽器本地。
  \item Constraint：零基礎開發，期一週內完成核心功能。
\end{itemize}

次步者，析其所求也。

\begin{enumerate}
  \item Step 1：實現「記錄學習條目」之功能（日期、內容、時長）；
  \item Step 2：增「學習時長統計」（今日、本週累計）；
  \item Step 3：實現「學習日曆」視圖（如 GitHub 貢獻圖之展示）；
  \item Step 4：增「學習回顧」之功能（隨機展示往日某天之學習記錄，以助複習）。
\end{enumerate}

第三步，撰一份簡短 PRD。不必長——上之拆解本身即構成 PRD 之主體。整理之為 .md 之檔，置於專案根目錄，以為君與 AI 之共同參考。

第四步，始第一步之開發。予 AI 之指令可如為：

\begin{lstlisting}
我須一個學習追蹤器之第一個功能——記錄學習條目。
用戶可添加一條學習記錄，包含日期、學習內容、學習時長（分鐘）。
數據存在localStorage。
前端用React + TypeScript。
請予出完整實現。
\end{lstlisting}

注意此指令之特點：惟述當前此步之目標（記錄功能），不摻他功能；明技術約束（React、TypeScript、localStorage）；予足夠信息使 AI 知「何所為」與「何以致之」，而具體實現由 AI 自決。此即經需求工程加工之後高品質之 AI 編程指令也。

\begin{tcbpractice}{來動手：為君之「心頭好」撰一份五要素卡}
擇一君醞釀已久之小專案，費 5 分鐘填一卡：Who / What / Why / How / Constraint，各一句話。乃拆之為 3-5 個「每一步皆能用」之步驟。君將見：念既書之於紙，即自「白日夢」化為「待辦清單」。

\end{tcbpractice}

\begin{tcbsummary}{本章小結}
\begin{itemize}
  \item 需求五要素：Who / What / Why / How / Constraint——思慮審明而後啟口
  \item 需求拆解黃金法則：每步既成，專案皆「能用之物」
  \item PRD 寫予自己：概述、功能列表、技術方案、資料模型、界面、計畫
  \item 善指令 = 惟述當前一步 + 明確約束 + 言「何所為」不言「何以致之」
  \item 完整推演模板：五要素 拆解 PRD 予指令，四步而行
\end{itemize}

\end{tcbsummary}
